\section*{Hoofdstuk \ref{ch:g1:materials-around-us} --- Waar dingen van gemaakt zijn}

\begin{solution}{exo:g1:materials-around-us:1}
Een stoel is een \omterm{def:g1:materials-around-us:object}{voorwerp}. Hout is een \omterm{def:g1:materials-around-us:material}{materiaal}: een stoel kan ervan
gemaakt zijn.
\end{solution}

\begin{solution}{exo:g1:materials-around-us:2}
Ruit: glas. Spijker: metaal. Kiezel: steen. Boek: papier. Sok: textiel
(wol of katoen).
\end{solution}

\begin{solution}{exo:g1:materials-around-us:3}
Hout: een potlood, een deur (of een stoel, een tafel). Plastic: een kom,
een emmer (of speelgoed, een fles).
\end{solution}

\begin{solution}{exo:g1:materials-around-us:4}
De metalen vork: dat is de enige die niet van hout of plastic is. (Twee
zijn van hout en één is van plastic.)
\end{solution}

\begin{solution}{exo:g1:materials-around-us:5}
Glas: het is doorzichtig.
\end{solution}

\begin{solution}{exo:g1:materials-around-us:6}
De stalen spijker. Een magneet trekt ijzer en staal aan, geen glas en
geen hout.
\end{solution}

\begin{solution}{exo:g1:materials-around-us:7}
Er drijven twee \omterm{def:g1:materials-around-us:object}{voorwerpen} (het houten blokje en de plastic dop); er
zinken er drie (de kiezel, de knikker en de spijker).
\end{solution}

\begin{solution}{exo:g1:materials-around-us:8}
De bladen zijn van metaal, de handvatten van plastic. Een schaar is van
twee materialen gemaakt.
\end{solution}

\begin{solution}{exo:g1:materials-around-us:9}
In de blauwe hoepel: het is van metaal. Het ligt niet in de oranje hoepel,
omdat een magneet het niet aantrekt: het is van aluminium, niet van
ijzer of staal.
\end{solution}

\begin{solution}{exo:g1:materials-around-us:10}
Een raam moet licht binnenlaten en we moeten erdoorheen kunnen kijken.
Glas is doorzichtig, hout niet.
\end{solution}

\begin{solution}{exo:g1:materials-around-us:11}
``Ja'': de spijker, de paperclip en de sleutel, 3 \omterm{def:g1:materials-around-us:object}{voorwerpen}. ``Nee'':
de kiezel, de lepel, het blikje, de knikker en de sok, 5 \omterm{def:g1:materials-around-us:object}{voorwerpen}.
$3 + 5 = 8$: elk \omterm{def:g1:materials-around-us:object}{voorwerp} heeft zijn plaats.
\end{solution}
